\section*{1. Introduction and Research Background}

The rapid diffusion of generative artificial intelligence is changing the basic way human beings acquire information, organize evidence, and form judgments under uncertainty. Unlike search engines, databases, and traditional recommender systems, generative AI does not merely help users find information more quickly, nor does it simply recommend one option to them. It is increasingly used to interpret highly uncertain issues such as disease risk, investment prospects, career choice, public policy, legal responsibility, historical controversy, and social conflict. In such settings, users are not typically facing a factual query with a clearly existing answer; rather, they are facing a judgment task in which the evidence is incomplete, the sources are heterogeneous, explanations compete with one another, and conclusions are probabilistic.

In a traditional information environment, individuals usually encounter evidence first and then gradually form explanations and judgments. Even though search engines and recommender systems reduce the cost of obtaining information, users still need to compare evidence across multiple sources, identify conflicts, assess reliability, and decide when to stop searching. In other words, traditional information technology mainly changes the efficiency and path of information search rather than fully replacing individuals' direct contact with the evidence space. Although users may be influenced by ranking, recommendation, and attention-allocation mechanisms, they are still usually aware that they face an information environment composed of multiple sources, multiple viewpoints, and multiple pieces of evidence.

Generative AI changes this epistemic order. Large language models usually do not first present users with a set of scattered information sources. Instead, they directly generate a structured, coherent, and seemingly well-synthesized explanation. As a result, users often encounter an algorithmically organized judgment framework before they have been fully exposed to the underlying evidence. This framework contains not only a conclusion, but also reasons, evidence ordering, conflict handling, and causal narrative. What users therefore face is no longer an unprocessed evidence environment, but an evidence environment that has already been pre-interpreted by an algorithm.

The key point of this change is that generative AI is not simply providing ``more information.'' It is completing in advance part of the epistemic work that users would otherwise have had to perform themselves. It helps users decide which evidence is more important, which conflicts can be interpreted as exceptions, which uncertainties can be compressed into understandable narratives, and which conclusions appear more reasonable under the current evidence. For users, the value of AI output comes not only from the amount of information it contains, but also from its formal completeness, fluency, and interpretability. A clearly structured, linguistically self-consistent, and logically coherent AI answer can give users the feeling that the main information has already been fully processed.

Accordingly, the social-scientific significance of generative AI should not be understood merely as a reduction in information-acquisition costs. More importantly, it may alter the basic process through which individuals form beliefs under uncertainty. In traditional belief formation, individuals usually confront the incompleteness, heterogeneity, and conflict of evidence first, and then gradually form a judgment within those uncertainties. In a generative AI environment, by contrast, individuals may first encounter a seemingly complete explanation and only then decide whether it is necessary to return to the evidence itself. In other words, cognitive closure may occur before users have adequately confronted evidentiary uncertainty.

From the perspective of behavioral economics and information economics, judgment under uncertainty contains at least two distinguishable stages. First, individuals decide whether to continue acquiring information. This stage involves the value of additional information, search costs, and optimal stopping. Second, individuals decide how to update beliefs after receiving new evidence. This stage involves evidentiary diagnosticity, prior beliefs, and Bayesian updating. The distinctive feature of generative AI is that it may intervene in both stages simultaneously. On the one hand, the AI's synthetic answer may lead users to believe mistakenly that the main evidence has already been sufficiently covered, thereby reducing the perceived value of continued search. On the other hand, the AI's coherent narrative may compress ambiguity and conflict within the problem, thereby reducing the degree to which users update on later disconfirming evidence.

This means that the risk of generative AI is not limited to whether it produces false information, biased content, or hallucinated facts. A subtler risk is that, even when the AI's answer is not obviously wrong, its highly fluent and structurally complete explanatory form may cause users to believe prematurely that they already understand the issue, or that the main evidence has already been adequately processed. The key questions then are no longer only ``Did the user adopt an incorrect answer?'' but also ``Did the user stop searching for evidence too early?'' and ``Did the user underestimate the importance of later disconfirming evidence?''

Existing research on algorithmic judgment mainly focuses on whether humans trust, adopt, or avoid algorithmic advice, such as algorithm aversion, algorithm appreciation, and automation bias \citep{ParasuramanRiley1997,LeeSee2004,Dietvorst2015,Logg2019}. This literature provides an important foundation for understanding human--machine decision-making, but it mostly studies traditional algorithmic advice, where the algorithm provides a predicted value, score, classification result, ranking result, or recommended option.

The difference between generative AI and traditional recommendation algorithms is not merely that the algorithm is more complex. More fundamentally, generative AI changes the form of information and the epistemic relationship users have to it. Traditional recommender systems mainly change the exposure structure of information through ranking, filtering, and personalized recommendation; their core function is to determine which discrete information items are more likely to enter the user's field of vision \citep{ResnickVarian1997,AdomaviciusTuzhilin2005,Bozdag2013}. Users may be influenced by rankings and recommendations, but they still typically confront multiple relatively separate information entry points and must compare, interpret, and integrate them on their own.

Generative AI further changes the representational form of information. It does not simply arrange external information items, but directly generates a continuous, fluent, and structurally complete natural-language explanation. Because of mechanisms such as the Transformer architecture, autoregressive language modeling, instruction tuning, and reinforcement learning from human feedback, the typical outputs of generative AI often appear contextually relevant, semantically coherent, and ``helpful'' \citep{Vaswani2017,Brown2020,Ouyang2022}. This means that generative AI folds evidence selection, evidence weighting, conflict handling, and conclusion generation into a single answer, so that before users have been fully exposed to the underlying evidence, they are already facing a knowledge product that appears to have completed the synthesis.

Thus, the influence of generative AI lies not only in changing which information users see, but also in changing how users understand \emph{how} the answer was generated. When facing a traditional recommendation algorithm, users are typically evaluating a ranking result. When facing generative AI, users are more likely to infer whether the AI has already understood the problem, seen enough evidence, covered the major dimensions, processed disconfirming evidence, and whether its summary represents the entire evidence environment. This is the theoretical starting point for the concept of \emph{perceived signal structure}.

To capture this mechanism, this proposal introduces the concept of perceived signal structure. Perceived signal structure refers to users' subjective beliefs about the information-generation process behind a generative AI summary. More concretely, users do not merely assess whether the AI's conclusion is credible; they also infer how much evidence the AI has observed, which dimensions of evidence it has covered, whether it has processed disconfirming information, whether it has integrated different sources, whether it has compressed evidentiary conflict, and whether the summary can represent a larger evidence pool. In other words, when users face AI output, they are not only receiving an answer; they are also forming a subjective model of the signal-generation process behind that answer.

This concept helps distinguish the present proposal from the existing literature on algorithmic trust. The influence of generative AI does not come entirely from ``algorithmic authority,'' nor does it arise solely from direct trust in the AI's conclusion. Even if users do not fully believe the AI's final judgment, they may still believe that the AI has already searched, filtered, and synthesized extensively. Even if users retain some doubts about the AI, they may still reduce the perceived necessity of continued search and the diagnosticity of later counterevidence because they believe the AI has already processed the main disconfirming evidence. The key influence of generative AI, therefore, is not simply to increase or decrease trust in the algorithm. Rather, it changes users' understanding of the current evidence environment itself.

From the perspective of information economics, this process can be understood as a special kind of information-design effect. The classic framework of Bayesian persuasion shows that an information sender can influence a receiver's posterior beliefs and action choices by designing the signal structure \citep{KamenicaGentzkow2011}. More broadly, the literature on information design studies how the designer of an information structure can shape receivers' optimal behavior under uncertainty \citep{BergemannMorris2019}. This proposal draws on that line of thought but shifts the focus from an external sender designing a signal structure to users' perceptions of the signal structure behind a generative AI summary.

Further, the literature on non-Bayesian persuasion shows that when receivers do not update according to the normative Bayesian rule, the effect of information design depends on the form of receivers' systematic inferential biases \citep{deClippelZhang2022}. This perspective is especially suitable for analyzing generative AI. AI does not merely transmit signals; in generating a summary, it reorganizes the signal structure. Users likewise do not merely update on a signal; before updating, they first form subjective inferences about the AI's signal-generation process. The influence of generative AI therefore depends not only on the actual signal structure, but also on the signal structure that users perceive.

On this basis, the proposal develops a framework of ``perceived signal structure and double Bayesian distortion.'' The core claim is that answer-first synthesis by generative AI may influence human belief formation under uncertainty through two mechanism paths.

First, at the information-acquisition stage, generative AI may induce an illusion of algorithmic exhaustiveness. Because AI summaries often appear in a structured, synthetic, and interpretive form, users may overestimate their evidentiary coverage, sample size, and representativeness. They may believe that the AI has already searched and integrated the major evidence sources, compared different viewpoints, and processed the key conflicts. As a result, the subjective value of additional search falls, the marginal return to continued information acquisition is underestimated, and individuals are more likely to stop evidence sampling too early.

Second, at the belief-updating stage, generative AI may produce an ambiguity-compression effect. The coherent narrative produced by AI tends to organize scattered, conflicting, and uncertain evidence into a clearer explanatory framework. When users believe that the AI has already processed the main disconfirming evidence, later counterevidence may be perceived as a signal with lower novelty, lower representativeness, or lower diagnosticity. As a result, users may encounter new disconfirming evidence without revising their prior beliefs to the degree warranted by its objective information content.

Generative AI may therefore create a \emph{double Bayesian distortion}. The first distortion occurs at the information-acquisition stage, where users underestimate the value of additional signals and stop searching too early. The second distortion occurs at the information-processing stage, where users underestimate the diagnosticity of disconfirming signals and update too little. The term ``double Bayesian distortion'' here does not refer only to posterior updating errors in a narrow sense. Rather, it refers to simultaneous deviations from the normative Bayesian decision process in both information acquisition and belief updating.

The central theoretical argument of this proposal is that the key influence of generative AI does not lie only in changing what answer users accept. By providing a coherent explanation before users see the evidence, it changes the rules by which humans acquire information and update beliefs under uncertainty. It may lead users to search less because they believe the evidence space has already been adequately covered, and it may lead them to update less because conflicting evidence has already been compressed into an existing narrative. In this way, generative AI can produce a double Bayesian distortion under conditions of algorithmic authority, simultaneously altering both the information-acquisition and information-processing stages.

\section*{2. Literature Review}

\subsection*{2.1 Information Acquisition and Evidence Sampling under Uncertainty}

Belief formation under uncertainty does not begin with the final judgment; it begins with the decision of whether to keep looking for evidence. What individuals ultimately believe depends not only on how they interpret the evidence they have already seen, but also on how much evidence they saw before the judgment was formed, which evidence they saw, and at what point they stopped searching. Information acquisition is not a neutral preparatory step prior to belief formation. It is the first key site at which belief distortion can arise.

From a normative perspective, information search can be understood as a cost--benefit trade-off. When the expected value of additional information is higher than the cost of obtaining it, continuing to search is rational. When the marginal value of additional information is lower than the search cost, stopping is also rational. The literatures on rational inattention and costly information acquisition both emphasize that information is not free; individuals' attention and search behavior should be understood as a process of information choice under cost constraints \citep{Sims2003,CaplinDean2015}.

Classic decision research likewise shows that people do not process all information under uncertainty costlessly and without limit. Instead, they adopt adaptive search strategies and decide when to stop searching based on task complexity, information cost, and expected return \citep{PayneBettmanJohnson1993}. In reality, however, stopping rules do not always conform to normative requirements. Individuals may seek supporting information because of confirmation bias, terminate search too early because of a need for cognitive closure, or reduce their willingness to continue exploring because the first information they receive appears too clear, satisfactory, or coherent \citep{Nickerson1998,LordRossLepper1979}.

Generative AI introduces a new question for research on information acquisition. When users see a highly structured, coherent, and seemingly well-synthesized AI summary before they encounter the underlying evidence, do they systematically overestimate the sufficiency of the current evidentiary state? Unlike a traditional search environment, an AI summary may make users feel that they have indirectly encountered a larger evidence space even though they have not actually seen the distribution, variance, and conflict structure of that evidence.

Accordingly, this proposal argues that AI may change users' stopping rules by changing their subjective imagination of the evidence space. Users do not search less simply because they are ``lazy.'' They search less because they infer that the AI has already completed a considerable amount of information search, filtering, and integration. This mechanism is incorporated into the framework of perceived signal structure: if users believe that an AI summary is based on more evidence, broader dimensions, and more representative evidence, then the subjective value of continued search will decline.

\subsection*{2.2 Belief Updating, Evidentiary Diagnosticity, and Under-Updating}

After acquiring information, individuals still need to form or revise beliefs on the basis of the evidence they have obtained. In the normative Bayesian framework, individuals should adjust posterior beliefs according to the diagnosticity of new evidence. If a piece of evidence substantially raises or lowers the probability that a hypothesis is true, the individual's posterior judgment should change accordingly.

Yet a large body of behavioral decision research shows that human belief updating frequently deviates from normative Bayesian rules. Individuals may display conservatism and adjust too little in the face of new evidence; they may be influenced by anchoring so that later judgments depend too strongly on initial information; or they may reduce the weight of disconfirming evidence through motivated reasoning or confirmation bias. These distortions suggest that belief updating depends not only on the objective diagnosticity of the evidence itself, but also on how individuals understand the relationship between that evidence and their preexisting explanatory framework.

This point is especially important for the present proposal. Disconfirming evidence does not automatically change people's beliefs. It must first be recognized as evidence that is relevant, credible, and diagnostically meaningful before it can generate substantial updating. If individuals interpret disconfirming evidence merely as an exception, noise, a boundary case, or information already absorbed by the original narrative, then even objectively highly diagnostic evidence may not be given sufficient weight.

Existing research on belief updating has paid more attention to the effects of priors, motivation, cognitive ability, or information presentation format, and less attention to how the explanatory framework generated in advance by generative AI changes users' perceptions of the diagnosticity of later evidence. This proposal argues that this is the second key gap in understanding how AI affects human judgment. Generative AI may influence not only whether users continue to search for evidence, but also how they interpret the same evidence after they see it.

This motivates the mechanism of \emph{ambiguity compression}. By providing a coherent, complete, and readable explanatory narrative, generative AI organizes evidence that was originally heterogeneous, conflicting, and probabilistic into a more stable framework of understanding. That framework may reduce users' perception of unresolved uncertainty and evidentiary conflict, making later disconfirming evidence seem less important, less representative, or less capable of changing the original judgment.

\subsection*{2.3 Algorithmic Advice, Algorithmic Trust, and Algorithmic Authority}

Existing research on how humans respond to algorithmic outputs mainly revolves around algorithmic advice, algorithmic trust, algorithm aversion, algorithm appreciation, and automation bias. This research shows that humans neither consistently trust nor consistently reject algorithms. Sometimes, after seeing an algorithm make a mistake, people overavoid algorithmic judgment. At other times, they regard algorithms as more objective and more accurate than humans and therefore become more willing to adopt algorithmic advice. Research on automation bias further shows that, in some contexts, humans may overrely on automated systems and thereby neglect their own judgment or external evidence.

This literature provides an important basis for understanding human--machine decision relations, but it usually treats algorithmic output as a prediction, score, classification result, or recommended option. As a consequence, the main concern is often whether users adopt the algorithmic advice, revise their own judgment, overrely on the algorithm, or become averse to it after the algorithm errs. In short, the existing literature primarily focuses on the effect of algorithmic output on final judgment or choice.

The distinctiveness of generative AI lies not only in the fact that it provides explanations, but also in the fact that users may understand it as an epistemic agent with a knowledge state and judgmental capacity. Traditional recommendation algorithms are usually understood as ranking tools, filters, or decision aids. By contrast, generative AI interacts with users through natural-language dialogue and can explain, ask follow-up questions, revise, summarize, and respond to specific contexts. This interactive format makes users more likely to understand it as a quasi-subject that has already ``seen,'' ``understood,'' and ``judged,'' rather than as a neutral computational tool.

Research in human--computer interaction has long shown that humans apply social rules and social expectations to computer systems. The CASA paradigm argues that when computers display language, responsiveness, interactivity, or social cues, people unconsciously treat them as social actors \citep{ReevesNass1996,NassMoon2000}. Research in social cognition likewise shows that people anthropomorphize and attribute minds to entities on the basis of interactivity, goal directedness, and behavioral complexity, inferring what they ``know,'' what they have ``understood,'' and whether they have ``already considered'' particular information \citep{EpleyWaytzCacioppo2007,Waytz2010}.

This point is especially important in the context of generative AI. Woodruff and Hewitt \citeyearpar{WoodruffHewitt2026} argue that, although LLMs can generate fluent and seemingly knowledgeable answers, their outputs do not satisfy knowledge conditions such as reflective regulation, resistance to revision, and normative commitment. They therefore should not be designed or understood as epistemic authority, but rather as knowledge-building partners \citep{WoodruffHewitt2026}. Qi and Pan \citeyearpar{QiPan2026} similarly argue in the context of evidence-based medicine that general-purpose LLMs can generate medically plausible evidence-related outputs, yet remain limited in source verifiability, methodological rigor, and responsibility attribution \citep{QiPan2026}.

Accordingly, the ``algorithmic authority'' discussed in this proposal is not just the general belief that ``the AI answer is correct.'' It refers, more specifically, to the user's tendency to regard AI as an epistemic agent that has already carried out the work of evidence search, evidence synthesis, and preliminary judgment on the user's behalf. When facing generative AI, users may ask not only ``Is this answer right?'' but also infer ``Has it already seen enough evidence?'', ``Has it already considered the opposing view?'', and ``Has it already completed the synthesis for me?'' These inferences are precisely what this proposal calls perceived signal structure.

The contribution of the present project to the literature on algorithmic advice is therefore as follows: existing studies mainly explain whether humans accept algorithmic output, whereas this proposal focuses on how algorithmic authority changes human judgments about evidentiary sufficiency and evidentiary diagnosticity.

\subsection*{2.4 Generative AI, Information Design, and Non-Bayesian Persuasion}

The literature on information design studies how a sender who controls the presentation of information can shape a receiver's beliefs and actions by designing the signal structure. In the classic Bayesian persuasion framework, the sender does not directly control the receiver's action but instead chooses an information structure that shapes the receiver's posterior beliefs and thereby influences action choice \citep{KamenicaGentzkow2011}. The later literature on information design extends this idea into a unified framework for analyzing the relationship among information structures, posterior beliefs, and behavioral responses \citep{BergemannMorris2019}. However, this traditional framework usually assumes that receivers can interpret signals correctly according to Bayesian rules. In reality, individuals often display systematic deviations in probabilistic inference. Non-Bayesian persuasion therefore emphasizes that when the receiver is not a Bayesian updater, the actual posterior can be understood as the result of applying some distortion function to the normative Bayesian posterior \citep{deClippelZhang2020}.

This proposal builds on that idea and extends it into the context of generative AI. Generative AI is not a simple extension of traditional recommendation algorithms. Search engines, recommender systems, and news-feed algorithms mainly change users' information-exposure structures through ranking, filtering, and personalization; their core function is to determine which discrete information items are more likely to enter the user's field of vision \citep{ResnickVarian1997,AdomaviciusTuzhilin2005,Bozdag2013}. In such systems, users are still usually facing multiple relatively separate information entry points, such as links, products, articles, videos, or ratings, and must compare, interpret, and integrate them on their own. In short, traditional recommendation algorithms mainly change what users ``see'' and what they ``see first.''

The crucial difference with generative AI is that it further changes both the expressive form of information and the epistemic relationship to it. Rather than arranging external information items, it directly generates a continuous, fluent, and structurally complete natural-language explanation. Because of mechanisms such as the Transformer architecture, autoregressive language modeling, instruction tuning, and reinforcement learning from human feedback, the typical outputs of generative AI often appear contextually relevant, semantically coherent, natural in tone, and ``helpful'' \citep{Vaswani2017,Brown2020,Ouyang2022}. What users face, therefore, is not a set of information sources waiting for further evaluation, but a synthetic answer that has already been organized, compressed, and narrativized by the model. Generative AI folds into a single output what were traditionally scattered across multiple cognitive stages: evidence selection, evidence weighting, conflict handling, causal explanation, and conclusion generation.

This gives generative AI a persuasive structure that differs from that of traditional recommendation algorithms. Traditional recommender systems may influence information exposure, but they still generally preserve visible separation among multiple information items. Generative AI instead integrates heterogeneous evidence, conflicting information, and uncertainty into a single, outwardly coherent explanatory narrative. Narrative transportation theory shows that coherent narratives can influence belief by absorbing attention, emotion, and imagination, thereby lowering the tendency to rebut or evaluate critically \citep{GreenBrock2000}. In the generative AI context, this mechanism may be amplified further: what AI outputs is not scattered evidence, but an explanatory framework that has already completed an initial act of meaning-making. Users may therefore more readily understand it as a stable posterior judgment rather than as a hypothesis that still requires testing.

Yet the formal coherence of generative AI is not equivalent to real understanding or adequate evidence synthesis. Large language models can generate highly fluent text, but linguistic fluency does not guarantee semantic grounding, factual sensitivity, or evidence-sensitive judgment. This is precisely the basis of the critique advanced by Bender et al. \citeyearpar{Bender2021}: language models operate mainly at the level of linguistic form, while humans easily infer meaning understanding and knowledge processing from fluent syntax and coherent discourse \citep{Bender2021}. The risk of generative AI, then, is not only that it may generate false content, but that users may misread its fluency and completeness as signs that the system has already adequately processed the underlying evidence.

More than that, users may understand generative AI as an epistemic agent with its own knowledge state and judgmental capacity. Traditional recommendation algorithms are usually viewed as back-end rankers, filters, or auxiliary tools. Generative AI, by contrast, responds in natural language, can explain, revise, summarize, acknowledge uncertainty, and adjust answers to context. Research in human--computer interaction shows that when computer systems display language, responsiveness, and social cues, humans tend to apply social rules to them and treat them as social actors \citep{ReevesNass1996,NassMoon2000}. Social cognition research likewise shows that people anthropomorphize and attribute mentality on the basis of interactivity, goal directedness, and behavioral complexity. In the generative AI context, these tendencies create fertile ground for users to infer more from the AI's answer than the answer itself contains.

\subsection*{2.5 Gaps in the Literature and the Positioning of This Proposal}

Taken together, the existing literature provides an important foundation for the present proposal, but it still leaves an insufficiently explained theoretical gap: existing studies have rarely analyzed how generative AI changes users' subjective understanding of \emph{how} signals are generated and, in turn, how that understanding affects information acquisition and belief updating.

First, the literature on information acquisition and evidence sampling shows that whether individuals continue searching depends on their judgments about the value of additional information, search cost, and the sufficiency of current evidence. Yet this literature generally focuses on users' own search costs, attentional limits, or motivational biases, and pays less attention to whether users overestimate the sufficiency of the current informational state when generative AI provides a synthetic explanation before they have contacted the original evidence. In other words, the existing literature has not adequately explained whether users reduce the subjective value of continued search because they believe the AI has already covered and integrated the main evidence.

Second, the literature on belief updating shows that people may display conservative updating, confirmation bias, or low weighting of disconfirming evidence, and that the subjective diagnosticity of evidence affects posterior revision. Yet this literature rarely examines how the coherent explanatory framework generated in advance by generative AI changes users' understanding of subsequent evidence. In particular, when users believe that the AI has already considered counterevidence, later disconfirming information may no longer be seen as new and highly diagnostic evidence, but as a local exception that has already been absorbed by the original narrative.

Third, the literature on algorithmic advice and algorithmic trust mainly focuses on whether users adopt, revise, or avoid algorithmic output, while continuing to interpret algorithmic output as a prediction, score, recommendation, or classification. The distinctive feature of generative AI is that it not only gives a conclusion, but also organizes evidence, explains causal relations, handles conflict, and generates coherent narratives in natural language. The key question is therefore not merely whether users believe the AI's answer, but how they understand the evidentiary basis of that answer: Has the AI seen enough evidence? Has it covered the main dimensions? Has it processed disconfirming information? Can its summary represent the overall evidence environment?

Fourth, the emerging literature on generative AI has already begun to address bias, hallucination, automation bias, narrative persuasion, and illusions of understanding, but it still needs to unpack the behavioral mechanisms behind those phenomena more directly. Existing research shows that AI outputs can affect judgment, but it rarely distinguishes between two different sources of deviation: does final judgment become distorted because users acquired less evidence earlier, or because they later assigned lower weight to the same evidence? This distinction is crucial both for understanding how AI affects belief formation and for designing appropriate governance interventions.

The present proposal is developed around this gap. It does not treat generative AI merely as a tool that provides advice. Rather, it treats generative AI as a meaning-making system that changes both users' signal environment and their understanding of signals. It introduces the concept of perceived signal structure to describe users' subjective inferences about the evidence-generation process behind an AI summary, including perceived sample size, coverage breadth, representativeness, degree of disconfirming-evidence processing, and compression loss. The argument is that answer-first synthesis by generative AI may affect both stages simultaneously: at the information-acquisition stage, users may overestimate the evidentiary coverage and representativeness of the AI summary and therefore reduce the subjective value of continued search; at the information-processing stage, users may believe that disconfirming evidence has already been considered or absorbed by the AI and therefore reduce the perceived diagnosticity of later counterevidence.

The question this proposal asks, then, is not whether AI affects human judgment in some general sense. It asks a more specific mechanism question: how does generative AI, by changing users' understanding of the signal-generation process, affect their judgments of evidentiary sufficiency and evidentiary diagnosticity under uncertainty, and how does that in turn produce insufficient evidence sampling and insufficient belief updating?
